%% ============================================================================
%% Companion PRD -- Gadallah & Liu.  REBUILT source-frame (PRD_REBUILD_PLAN).
%% Every value via % prd_numbers.tex -- GENERATED by prd_build_data.py; do NOT hand-edit.
% Source frame, z=0.0856. Every headline number of the companion PRD.
\def\Minf{62.71}
\def\chiinf{0.676}
\def\Ainf{171.7}
\def\Apre{108.5}
\def\Mtwelve{65.83}
\def\zval{0.0856}
\def\Mearly{75.4}
\def\MearlyDet{81.8}
\def\dMzero{12.6}
\def\dMsig{5.2}
\def\chiEarly{0.764}
\def\chiLate{0.49}
\def\lateFiveMean{57.9}
\def\Mfinal{56.2}
\def\AearlySrc{2.35}
\def\AlateSrc{1.48}
\def\Rareamin{1.31}
\def\RareaminEp{20}
\def\minPA{0.94}
\def\crossTime{4.4}
\def\mmEarly{67.5}
\def\mmEarlyDet{73.3}
\def\nrEarly{76.7}
\def\nrLate{56.4}
\def\nrTrue{63.0}
\def\rhoM{-0.974}
\def\rhoChi{-0.813}
\def\rhoA{-0.988}
\def\pblockM{0.008}
\def\pblockA{0.008}
\def\pthinM{0.017}
\def\pthinChi{0.083}
\def\tauU{1.9}
\def\tauW{2.6}
\def\tauEarly{4.0}
\def\redChiTwo{0.46}
\def\areaRatio{1.581}
\def\areaPct{58.1}
\def\areaDraws{0/20000}
\def\rhoIsaacson{0.81}
 (generated by code/prd_build_data.py,
%% gated). NO hand-typed numbers. Source frame; detector only in "(...detector)".
%% Honesty rules (Part F): no "clean convergence"; no "tau=tau220" as a
%% measurement; reduced chi^2 is the sigma-normalised value.
%% ============================================================================
\documentclass[aps,prd,twocolumn,superscriptaddress,nofootinbib,longbibliography]{revtex4-2}
\usepackage{amsmath,amssymb,graphicx}
\usepackage[colorlinks=true,linkcolor=blue,citecolor=blue,urlcolor=blue]{hyperref}
% prd_numbers.tex -- GENERATED by prd_build_data.py; do NOT hand-edit.
% Source frame, z=0.0856. Every headline number of the companion PRD.
\def\Minf{62.71}
\def\chiinf{0.676}
\def\Ainf{171.7}
\def\Apre{108.5}
\def\Mtwelve{65.83}
\def\zval{0.0856}
\def\Mearly{75.4}
\def\MearlyDet{81.8}
\def\dMzero{12.6}
\def\dMsig{5.2}
\def\chiEarly{0.764}
\def\chiLate{0.49}
\def\lateFiveMean{57.9}
\def\Mfinal{56.2}
\def\AearlySrc{2.35}
\def\AlateSrc{1.48}
\def\Rareamin{1.31}
\def\RareaminEp{20}
\def\minPA{0.94}
\def\crossTime{4.4}
\def\mmEarly{67.5}
\def\mmEarlyDet{73.3}
\def\nrEarly{76.7}
\def\nrLate{56.4}
\def\nrTrue{63.0}
\def\rhoM{-0.974}
\def\rhoChi{-0.813}
\def\rhoA{-0.988}
\def\pblockM{0.008}
\def\pblockA{0.008}
\def\pthinM{0.017}
\def\pthinChi{0.083}
\def\tauU{1.9}
\def\tauW{2.6}
\def\tauEarly{4.0}
\def\redChiTwo{0.46}
\def\areaRatio{1.581}
\def\areaPct{58.1}
\def\areaDraws{0/20000}
\def\rhoIsaacson{0.81}
   % generated \def macros -- do not hand-edit numbers

\newcommand{\Meff}{M_{\rm eff}}\newcommand{\aeff}{\chi_{\rm eff}}
\newcommand{\Aeff}{A_{\rm eff}}\newcommand{\Msun}{M_\odot}\newcommand{\tstart}{t_0}

\begin{document}
\title{Quantifying Single-Mode Kerr Inference Bias Through\\
Time-Resolved Ringdown Analysis of GW250114}
\author{Attia A. Gadallah}
\author{Zi-Kui Liu}
\affiliation{Department of Materials Science and Engineering, The Pennsylvania
State University, University Park, Pennsylvania 16802, USA}
\date{\today}

\begin{abstract}
Fitting a black-hole ringdown with a single quasi-normal mode (QNM) started too
early returns biased remnant parameters. We quantify this bias for GW250114 by
sliding the analysis start time $\tstart$ across 21 overlapping epochs and
tracking the inferred Kerr mass, spin and horizon area. The single-mode
effective mass begins at $\Mearly\,\Msun$ (source frame)---above the total
pre-merger mass $M_1{+}M_2=\Mtwelve\,\Msun$, a physically impossible value for
any true remnant---and decreases monotonically ($\rho_s=\rhoM$, block-permutation
$p=\pblockM$), crossing the global posterior anchor $M_\infty=\Minf\,\Msun$ near
$\crossTime$~ms and undershooting it at late times as the signal-to-noise
degrades. A controlled numerical-relativity injection reproduces the same
overshoot-to-undershoot behaviour about its \emph{known} true mass, confirming
the trajectory as inference bias rather than physical evolution. Three
independent tests attribute the early excess to overtone spectral content, and
the inferred horizon area exceeds the progenitor sum at every epoch
[$\Aeff/(A_1{+}A_2)\ge\Rareamin$], consistent with the Hawking area theorem
($A_\infty/(A_1{+}A_2)=\areaRatio$, $+\areaPct\%$, $\areaDraws$ draws). Because
the inference chain operates in the detector frame, we quote all quantities in
the source frame throughout; the relaxation timescale is window-limited and not
a clean measurement of $\tau_{220}$. A companion Letter recasts these
systematics as an exact balance law for the projected horizon.
\end{abstract}
\maketitle

\section{Introduction}
Gravitational-wave ringdown spectroscopy infers a remnant black hole's mass and
spin from the QNM spectrum of the post-merger signal
\cite{Berti2006,Giesler2019,Isi2019}. Three distinct masses enter: the total
pre-merger mass $M_1{+}M_2=\Mtwelve\,\Msun$; the true remnant mass
$M_\infty=\Minf\,\Msun$ from the full inspiral--merger--ringdown posterior; and
the \emph{inferred} effective mass $\Meff(\tstart)$ returned by a single-mode
fit begun at start time $\tstart$. When $\tstart$ is early, unmodeled overtones
and non-linear transients inflate $\Meff$ above $M_\infty$
\cite{Giesler2019,Isi2019,Bhagwat2018}; as $\tstart$ advances and the overtones
decay, the fit relaxes toward the fundamental mode. This work maps that
relaxation quantitatively for GW250114 \cite{LVK_GW250114}, and shows that its
early value is not merely biased but \emph{physically impossible}, exceeding the
total available pre-merger mass. A companion Letter \cite{Companion2026} recasts
the same trajectory as an exact first law on the manifold of Kerr equilibria; the
present paper is the observational study.

\section{Data and methods}
We use the $2\times10^4$ Bilby NRSur7dq4 posterior samples of the LVK analysis
\cite{Ashton2019,Varma2019} for the global anchors, and the per-epoch
$(f_{220},\tau_{220})$ posteriors of the $220{+}440$ pyRing analysis
\cite{Carullo2019,Cornish2015} for the single-mode trajectory. Windows span
$4.59$~ms (detector), stepped by $1M=0.337$~ms, giving $92.7\%$ overlap and 21
epochs to $6.7$~ms. Within each epoch we invert $(f,\tau)\to(M,\chi)$ per sample
through the Berti fitting relations \cite{Berti2006},
$\chi$ from $Q=\pi f\tau$ and $M$ from $f$, then form $A=8\pi M^2(1+\sqrt{1-
\chi^2})$.

\emph{Frame.}---The pyRing chain fits the observed (detector-frame) strain and
returns detector-frame masses (its configuration remnant reference is
$68.41\approx M_\infty(1{+}z)$, and the late observed $f_{220}$ matches the
detector-frame remnant frequency). We therefore convert every quantity to the
source frame by the single Bilby median factor $1{+}z$ ($z=\zval$):
$M_{\rm src}=M_{\rm det}/(1{+}z)$, $A_{\rm src}=A_{\rm det}/(1{+}z)^2$, with
$\chi$ frame-invariant. All values below are source-frame; detector values
appear only in explicit ``($\ldots$ detector)'' parentheticals. This conversion
is essential: comparing the detector trajectory (early $\MearlyDet$~detector)
against the source anchor $\Minf$ inflates the excess and manufactures a false
convergence.

\section{Results: the inferred-parameter trajectory}
Figure~\ref{fig:traj} shows $\Meff(\tstart)$ and the area ratio
$\Aeff/(A_1{+}A_2)$. The single-mode mass begins at $\Meff(0)=\Mearly\,\Msun$
(detector $\MearlyDet$), an excess of $\dMzero\,\Msun$ ($\dMsig\sigma$) above the
anchor. Crucially this early value \emph{exceeds the total pre-merger mass}
$M_1{+}M_2=\Mtwelve\,\Msun$ by $14.4\%$: no true remnant can be heavier than the
matter that formed it, so the early single-mode mass is model-independently
unphysical---the cleanest demonstration that the trajectory is inference bias.
The mass then decreases monotonically (Spearman $\rho_s=\rhoM$;
block-permutation $p=\pblockM$; thinned exact-permutation $p=\pthinM$), the spin
from $\chiEarly$ to $\chiLate$, and the area monotonically likewise
($\rho_s=\rhoA$).

\begin{figure}[t]
\includegraphics[width=\columnwidth]{figures/fig02_trajectory.pdf}
\caption{\label{fig:traj}Single-mode inferred trajectory for GW250114 (source
frame). (a)~$\Meff(\tstart)$ (90\% band) starts at $\Mearly\,\Msun$ inside the
physically forbidden region $M>M_1{+}M_2$ (shaded), crosses the Bilby anchor
$M_\infty$ near $\crossTime$~ms, and undershoots at late times.
(b)~The horizon-area ratio $\Aeff/(A_1{+}A_2)$ exceeds unity at every epoch
(minimum $\Rareamin$), remaining above the Hawking bound.}
\end{figure}

\section{Late-epoch behaviour}
The trajectory does not converge cleanly to the anchor: it \emph{crosses}
$M_\infty$ near $\crossTime$~ms and undershoots, the mean of the last five
epochs being $\lateFiveMean\,\Msun$ and the final epoch $\Mfinal\,\Msun$. This
undershoot is the expected low-signal-to-noise behaviour of the QNM inversion:
at late times the damping time is poorly constrained, biasing the quality factor
and hence $\chi$ (and through it $M$) downward. The picture is therefore
overshoot ($\tstart\lesssim\crossTime$~ms) $\to$ crossing $\to$
noise-dominated undershoot, not monotonic convergence. The spin channel is the
most precision-limited, its credible interval widening by a large factor from
early to late epochs (Fig.~\ref{fig:late}); we do not read a physical trend
from its late values.

\begin{figure}[t]
\includegraphics[width=\columnwidth]{figures/fig03_late_epoch.pdf}
\caption{\label{fig:late}Late-epoch behaviour (source frame).
(a)~Tension $(X-X_\infty)/\sigma_X$ for the mass, spin and area channels:
each overshoots the anchor at early $\tstart$, crosses near $\crossTime$~ms,
and undershoots at late times. (b)~The spin credible-interval width grows
sharply toward late epochs, the precision degradation that drives the
undershoot.}
\end{figure}

\section{Relaxation timescale}
A single-exponential fit to the mass excess is an adequate summary statistic
(reduced $\chi^2=\redChiTwo$, using the per-epoch credible-interval widths as
uncertainties), but the fitted timescale is \emph{not} a clean measurement of
the fundamental QNM damping time: across weighting and window choices it spans
$\tauU$--$\tauEarly$~ms, and the exponential model breaks down once the excess
crosses zero. The apparent timescale is window-limited---the analysis window
$T_w=4.59$~ms is itself of order $\tau_{220}$---so we report it only as
``order of $\tau_{220}$, window-limited'' and defer a full treatment to the
companion Letter, which shows the fit-response scale and $\tau_{220}$ are
structurally degenerate for one-cycle windows.

\section{Controlled numerical-relativity injection}
To confirm the trajectory is inference bias rather than physical evolution, we
apply the identical pipeline to a numerical-relativity injection with a
\emph{known} true remnant mass $\nrTrue\,\Msun$ (source; $68.41$ detector).
Figure~\ref{fig:nr} shows the single-mode recovery: it begins at
$\nrEarly\,\Msun$, well above the injected truth, and undershoots to
$\nrLate\,\Msun$ late---the same overshoot-to-undershoot pattern seen in the
real event, now against a mass we control. Because the true value is known, this
rules out a physical interpretation of the descent and validates the windowing
pipeline.

\begin{figure}[t]
\includegraphics[width=\columnwidth]{figures/fig05_nr_control.pdf}
\caption{\label{fig:nr}Controlled NR injection. The single-mode recovery
(source frame) overshoots the known injected true mass ($\nrTrue\,\Msun$,
dashed) at early start times and undershoots it late, mirroring the real event
and establishing the trajectory as inference bias.}
\end{figure}

\section{Attribution of the early excess to overtones}
A single-mode Kerr signal is strongly favoured over noise at every start time
(the Bayes factor $\ln\mathcal{B}$ declines from $\sim\!485$ at $\tstart=0$ to
$\sim\!30$ at late times as signal is discarded), while independent
overtone-detection statistics are large only at the earliest epochs
(Fig.~\ref{fig:rob})---the overtone content is confined to the early ringdown,
exactly where the mass excess is largest.

\begin{figure}[t]
\includegraphics[width=\columnwidth]{figures/fig01_robustness.pdf}
\caption{\label{fig:rob}Ringdown robustness (source frame).
(a)~The single-mode signal Bayes factor $\ln\mathcal{B}_{\rm signal}$ declines
with start time but remains large throughout: the signal is favoured at all
epochs. (b)~The QNMRF overtone-detection statistic ($221{:}220$) is significant
only at the earliest start times, confining overtone content to the early
ringdown.}
\end{figure}

Three independent, data-driven tests tie the early-epoch excess to overtone
spectral content. (i)~\emph{Isaacson energy balance}: the start-time decay rate
of the mass excess tracks the Isaacson luminosity of the non-fundamental content
reconstructed from the multi-mode ($220{+}221{+}222$) posteriors, Spearman
$\rho_s=\rhoIsaacson$. (ii)~\emph{Multi-mode accounting}: fitting the same
windows with a multi-mode template returns a mass consistent with the anchor,
the early multi-mode value being $\mmEarly\,\Msun$ (detector $\mmEarlyDet$) with
a residual below $2.2\sigma$ at every epoch. (iii)~\emph{Spectral correlation}:
the single-mode excess correlates with the overtone spectral energy across the
common epochs. The three lines of evidence (Fig.~\ref{fig:attr}), each frame-robust
(correlations and ratios), consistently attribute the excess to unmodeled
overtones absorbed by the truncated template.

\begin{figure*}[t]
\includegraphics[width=\textwidth]{figures/fig06_attribution.pdf}
\caption{\label{fig:attr}The three overtone-attribution tests (source frame,
points coloured by start time). (a)~T1: the mass-excess decay rate
$-d\Meff/d\tstart$ against the residual Isaacson luminosity
($\rho_s=\rhoIsaacson$). (b)~T2: the multi-mode inferred mass lies near the
Bilby anchor (dashed) at early epochs, residual $\le2.2\sigma$. (c)~T3: the
single-mode excess $\Delta M$ against the overtone spectral energy
$E_{\rm ovt}$ ($\rho_s\simeq0.8$).}
\end{figure*}

A time-resolved Bayesian model comparison corroborates this: the per-epoch
expected log pointwise predictive density favours a two-mode ($220{+}221$) over
a single-mode template only at the earliest epochs, the preference vanishing at
a crossover near $2.8$~ms (Fig.~\ref{fig:waic})---precisely where the
overtone-detection statistic and the mass excess also fade.

\begin{figure}[t]
\includegraphics[width=\columnwidth]{figures/fig07_waic.pdf}
\caption{\label{fig:waic}Per-epoch model comparison (WAIC): the two-mode
$-$ single-mode expected-log-density difference is positive only at early start
times, with a crossover near $2.8$~ms (source frame), independently locating the
extent of resolvable overtone content.}
\end{figure}

\section{Horizon area and the Hawking theorem}
Although the inferred single-mode area decreases with start time (as the
overtone-inflated early value relaxes), it never falls below the progenitor
sum: $\Aeff(\tstart)\ge A_1{+}A_2$ at every epoch, with minimum ratio
$\Rareamin$ (epoch $\RareaminEp$) and sample-level probability
$P(\Aeff>A_1{+}A_2)\ge\minPA$. The final remnant area from the full posterior
exceeds the progenitor sum by $\areaPct\%$
($A_\infty/(A_1{+}A_2)=\areaRatio$, positive in $\areaDraws$ Monte-Carlo draws,
a one-sided Clopper--Pearson significance $\gtrsim3.6\sigma$; Fig.~\ref{fig:area}),
consistent with the Hawking area theorem and extending the GW150914 test
\cite{Isi2021} to this event. The monotonic decrease of the \emph{inferred}
area is not a violation: it reflects the convergence of the truncated
single-mode description toward the stationary Kerr geometry, not physical
horizon dynamics.

\begin{figure}[t]
\includegraphics[width=\columnwidth]{figures/fig04_area_theorem.pdf}
\caption{\label{fig:area}Hawking area theorem for GW250114: the posterior
distribution of $R=A_\infty/(A_1{+}A_2)$ lies entirely above unity (median
$\areaRatio$, $\areaDraws$ draws), a $+\areaPct\%$ irreversible area increase.}
\end{figure}

\section{Limitations and conclusion}
The 21 epochs overlap and are not independent; the earliest windows are the most
sensitive to overtone content and non-linear contamination; the results pertain
to a single event. Within these limits, the corrected, source-frame analysis
shows that a single-mode ringdown fit of GW250114 returns a start-time-dependent
Kerr trajectory whose early value is \emph{physically impossible}
($\Mearly>M_1{+}M_2$), that relaxes through the true remnant and undershoots
under noise, that is reproduced in a controlled NR injection about a known
truth, and whose excess is attributable to overtone content by three independent
tests---all while the horizon area respects the Hawking bound. The relaxation
timescale is window-limited and not a measurement of $\tau_{220}$. The central
conclusions---the physical-impossibility clincher, the overtone attribution, the
area bound---are built from frame-invariant inequalities, ratios and
correlations, and are unchanged by the detector/source frame convention that
governs the absolute scale. The companion Letter \cite{Companion2026} develops
the underlying mechanics and thermodynamics of this projected horizon.

\begin{acknowledgments}
This research used data and software from the Gravitational Wave Open Science
Center (\url{https://gwosc.org}), a service of LIGO Laboratory, the LIGO
Scientific Collaboration, Virgo and KAGRA. We thank the developers of
\textsc{Bilby}, \textsc{pyRing}, \textsc{qnm} and the NR surrogate models. Work
conducted at the PhasesLab, Penn State.
\end{acknowledgments}

\appendix
\section{Jensen-inequality area audit}
Because the Kerr area is non-linear in $(M,\chi)$, computing it from the
posterior medians can bias it relative to the sample-median area. A
sample-by-sample audit over the per-epoch posteriors gives a relative bias of
mean $-0.03\%$ (range $[-0.17,+0.12]\%$), confirming that median-based area
point estimates are unbiased at the precision of this analysis.

\begin{thebibliography}{99}
\bibitem{Berti2006} E.~Berti, V.~Cardoso, C.~M. Will, Phys. Rev. D \textbf{73}, 064030 (2006).
\bibitem{Giesler2019} M.~Giesler, M.~Isi, M.~A. Scheel, S.~A. Teukolsky, Phys. Rev. X \textbf{9}, 041060 (2019).
\bibitem{Isi2019} M.~Isi \emph{et al.}, Phys. Rev. Lett. \textbf{123}, 111102 (2019).
\bibitem{Bhagwat2018} S.~Bhagwat \emph{et al.}, Phys. Rev. D \textbf{97}, 104065 (2018).
\bibitem{Isi2021} M.~Isi \emph{et al.}, Phys. Rev. Lett. \textbf{127}, 011103 (2021).
\bibitem{LVK_GW250114} LIGO Scientific, Virgo, and KAGRA Collaborations, ``Black Hole Spectroscopy and Tests of General Relativity with GW250114,'' arXiv:2509.08099 [gr-qc] (2025), LIGO-P2500461; data release: Zenodo, \href{https://doi.org/10.5281/zenodo.17018009}{doi:10.5281/zenodo.17018009}.
\bibitem{Companion2026} A.~A. Gadallah, Z.-K. Liu, ``Mechanics and thermodynamics of projected horizons in black-hole ringdown,'' companion Letter (2026).
\bibitem{Carullo2019} G.~Carullo, W.~Del Pozzo, J.~Veitch, Phys. Rev. D \textbf{99}, 123029 (2019).
\bibitem{Ashton2019} G.~Ashton \emph{et al.}, Astrophys. J. Suppl. \textbf{241}, 27 (2019).
\bibitem{Varma2019} V.~Varma \emph{et al.}, Phys. Rev. Research \textbf{1}, 033015 (2019).
\bibitem{Cornish2015} N.~J. Cornish, T.~B. Littenberg, Class. Quantum Grav. \textbf{32}, 135012 (2015).
\end{thebibliography}
\end{document}
